% GENERATED FILE. Edit canonical lessons, then run npm run generate:books.
% canonical-source-hash: fnv1a64:cded39a2bb2c6609
% canonical-lessons: SA-C357-samitih, SA-C357-parisad, SA-C357-pramukhah, SA-C357-sacivah, SA-C357-nirdisati

\chapter{A Committee Meets}
\label{ch:sa-a-committee-meets}
\hlchaptermodality{357}

\begin{chapteropening}
\textbf{By the end of this chapter:} \emph{I can say who heads a committee and what it decides.}

Complete the last lesson of chapter 357: I can say who heads a committee and what it decides.
\end{chapteropening}

\section[\texorpdfstring{\sk{समितिः} (samitiḥ)}{समितिः (samitiḥ)}]{\sk{समितिः} (samitiḥ) — a committee}

\label{lesson:SA-C357-samitih}



\begin{quote}
Before the new one: say the Sanskrit for a department, then the Sanskrit for a proposal.
\end{quote}

\subsection*{You'll want to know: \sk{समितिः}}
\textbf{\sk{समितिः}} — \emph{samitiḥ} — \textquotedblleft{}a committee\textquotedblright{}.

\begin{quote}
\sk{समितिः} \sk{श्वः} \sk{निर्णयं} \sk{करिष्यति।}
\end{quote}

\begin{quote}
\emph{samitiḥ śvaḥ nirṇayaṃ kariṣyati.}
\end{quote}

\begin{quote}
The committee will decide tomorrow.
\end{quote}

A feminine noun.

\begin{cousinweb}
\textbf{\sk{समितिः}} (\emph{samitiḥ}) is built on \emph{sam}, \textquotedblleft{}together\textquotedblright{}, and the root √i, \textquotedblleft{}to go\textquotedblright{}; MW gives \emph{sam-iti} \textquotedblleft{}a meeting, an assembly, a council\textquotedblright{}, and \textquotedblleft{}a committee\textquotedblright{} is its modern sense (MW).
\end{cousinweb}

\begin{grammarlens}[title={What you've built}]
One more word for giving reasons and opinions.
\end{grammarlens}

\subsection*{Your turn}
\begin{itemize}
\raggedright
  \item \emph{Say it:} \emph{samitiḥ}
  \item \emph{Say it:} \emph{samitiḥ}, once more
  \item \emph{Say it:} \emph{samitiḥ śvaḥ nirṇayaṃ kariṣyati.} — The committee will decide tomorrow.
  \item \emph{Recall:} say the Sanskrit for a department, then the Sanskrit for a proposal, then say \emph{samitiḥ} again
\end{itemize}

\begin{tcolorbox}[breakable,colback=teal!4,colframe=teal!35!black,title={Before you move on}]
What does \sk{समितिः} mean? (\textquotedblleft{}A committee\textquotedblright{}.) Say it once more.
\end{tcolorbox}
\section[\texorpdfstring{\sk{परिषद्} (pariṣad)}{परिषद् (pariṣad)}]{\sk{परिषद्} (pariṣad) — a council}

\label{lesson:SA-C357-parisad}



\begin{quote}
Before the new one: say the Sanskrit for a proposal, then the Sanskrit for a committee.
\end{quote}

\subsection*{You'll want to know: \sk{परिषद्}}
\textbf{\sk{परिषद्}} — \emph{pariṣad} — \textquotedblleft{}a council\textquotedblright{}.

\begin{quote}
\sk{ग्रामस्य} \sk{परिषत्} \sk{जलस्य} \sk{व्यवस्थां} \sk{करोति।}
\end{quote}

\begin{quote}
\emph{grāmasya pariṣat jalasya vyavasthāṃ karoti.}
\end{quote}

\begin{quote}
The village council arranges the water supply.
\end{quote}

A feminine noun.

The stem ends in \emph{-d}; said on its own, as here, the word ends in \emph{-t}: \sk{परिषत्}, \emph{pariṣat}.

\begin{cousinweb}
\textbf{\sk{परिषद्}} (\emph{pariṣad}) is built on \emph{pari}, \textquotedblleft{}around\textquotedblright{}, and the root √sad, \textquotedblleft{}to sit\textquotedblright{}; MW gives \emph{pari-ṣad} \textquotedblleft{}an assembly, a meeting, an audience, a council\textquotedblright{} (MW).
\end{cousinweb}

\begin{grammarlens}[title={What you've built}]
One more word for giving reasons and opinions.
\end{grammarlens}

\subsection*{Your turn}
\begin{itemize}
\raggedright
  \item \emph{Say it:} \emph{pariṣad}
  \item \emph{Say it:} \emph{pariṣad}, once more
  \item \emph{Say it:} \emph{grāmasya pariṣat jalasya vyavasthāṃ karoti.} — The village council arranges the water supply.
  \item \emph{Recall:} say the Sanskrit for a proposal, then the Sanskrit for a committee, then say \emph{pariṣad} again
\end{itemize}

\begin{tcolorbox}[breakable,colback=teal!4,colframe=teal!35!black,title={Before you move on}]
What does \sk{परिषद्} mean? (\textquotedblleft{}A council\textquotedblright{}.) Say it once more.
\end{tcolorbox}
\section[\texorpdfstring{\sk{प्रमुखः} (pramukhaḥ)}{प्रमुखः (pramukhaḥ)}]{\sk{प्रमुखः} (pramukhaḥ) — foremost; chief}

\label{lesson:SA-C357-pramukhah}



\begin{quote}
Before the new one: say the Sanskrit for a committee, then the Sanskrit for a council.
\end{quote}

\subsection*{You'll want to know: \sk{प्रमुखः}}
\textbf{\sk{प्रमुखः}} — \emph{pramukhaḥ} — \textquotedblleft{}foremost; chief\textquotedblright{}.

\begin{quote}
\sk{सः} \sk{सभायाः} \sk{प्रमुखः} \sk{सदस्यः} \sk{अस्ति।}
\end{quote}

\begin{quote}
\emph{saḥ sabhāyāḥ pramukhaḥ sadasyaḥ asti.}
\end{quote}

\begin{quote}
He is the foremost member of the assembly.
\end{quote}

\begin{cousinweb}
\textbf{\sk{प्रमुखः}} (\emph{pramukhaḥ}) is built from \emph{pra}, \textquotedblleft{}before\textquotedblright{}, and \emph{mukha}, \textquotedblleft{}the face\textquotedblright{}; MW gives \emph{pra-mukha} \textquotedblleft{}first, foremost, chief, principal\textquotedblright{} (MW).
\end{cousinweb}

\begin{grammarlens}[title={What you've built}]
One more word for giving reasons and opinions.
\end{grammarlens}

\subsection*{Your turn}
\begin{itemize}
\raggedright
  \item \emph{Say it:} \emph{pramukhaḥ}
  \item \emph{Say it:} \emph{pramukhaḥ}, once more
  \item \emph{Say it:} \emph{saḥ sabhāyāḥ pramukhaḥ sadasyaḥ asti.} — He is the foremost member of the assembly.
  \item \emph{Recall:} say the Sanskrit for a committee, then the Sanskrit for a council, then say \emph{pramukhaḥ} again
\end{itemize}

\begin{tcolorbox}[breakable,colback=teal!4,colframe=teal!35!black,title={Before you move on}]
What does \sk{प्रमुखः} mean? (\textquotedblleft{}Foremost; chief\textquotedblright{}.) Say it once more.
\end{tcolorbox}
\section[\texorpdfstring{\sk{सचिवः} (sacivaḥ)}{सचिवः (sacivaḥ)}]{\sk{सचिवः} (sacivaḥ) — a secretary; a counsellor}

\label{lesson:SA-C357-sacivah}



\begin{quote}
Before the new one: say the Sanskrit for a council, then the Sanskrit for foremost.
\end{quote}

\subsection*{You'll want to know: \sk{सचिवः}}
\textbf{\sk{सचिवः}} — \emph{sacivaḥ} — \textquotedblleft{}a secretary; a counsellor\textquotedblright{}.

\begin{quote}
\sk{सचिवः} \sk{सभायाः} \sk{कार्यं} \sk{करोति।}
\end{quote}

\begin{quote}
\emph{sacivaḥ sabhāyāḥ kāryaṃ karoti.}
\end{quote}

\begin{quote}
The secretary does the assembly's work.
\end{quote}

A masculine noun.

\begin{cousinweb}
\textbf{\sk{सचिवः}} (\emph{sacivaḥ}) MW gives \emph{saciva} \textquotedblleft{}an associate, a companion\textquotedblright{} and \textquotedblleft{}a king's counsellor, a minister\textquotedblright{}, and \textquotedblleft{}a secretary\textquotedblright{} is its modern sense; MW gives no derivation, so its origin is uncertain (MW).
\end{cousinweb}

\begin{grammarlens}[title={What you've built}]
One more word for giving reasons and opinions.
\end{grammarlens}

\subsection*{Your turn}
\begin{itemize}
\raggedright
  \item \emph{Say it:} \emph{sacivaḥ}
  \item \emph{Say it:} \emph{sacivaḥ}, once more
  \item \emph{Say it:} \emph{sacivaḥ sabhāyāḥ kāryaṃ karoti.} — The secretary does the assembly's work.
  \item \emph{Recall:} say the Sanskrit for a council, then the Sanskrit for foremost, then say \emph{sacivaḥ} again
\end{itemize}

\begin{tcolorbox}[breakable,colback=teal!4,colframe=teal!35!black,title={Before you move on}]
What does \sk{सचिवः} mean? (\textquotedblleft{}A secretary; a counsellor\textquotedblright{}.) Say it once more.
\end{tcolorbox}
\section[\texorpdfstring{\sk{निर्दिशति} (nirdiśati)}{निर्दिशति (nirdiśati)}]{\sk{निर्दिशति} (nirdiśati) — directs; instructs}

\label{lesson:SA-C357-nirdisati}



\begin{quote}
Before the new one: say the Sanskrit for foremost, then the Sanskrit for a secretary.
\end{quote}

\subsection*{You'll want to know: \sk{निर्दिशति}}
\textbf{\sk{निर्दिशति}} — \emph{nirdiśati} — \textquotedblleft{}directs; instructs\textquotedblright{}.

\begin{quote}
\sk{अध्यक्षः} \sk{सदस्यान्} \sk{निर्दिशति।}
\end{quote}

\begin{quote}
\emph{adhyakṣaḥ sadasyān nirdiśati.}
\end{quote}

\begin{quote}
The president directs the members.
\end{quote}

Past (he …): \sk{सः} \sk{सदस्यान्} \sk{निर्दिष्टवान्}, \emph{saḥ sadasyān nirdiṣṭavān}.

\begin{cousinweb}
\textbf{\sk{निर्दिशति}} (\emph{nirdiśati}) is built on \emph{nir}, \textquotedblleft{}out\textquotedblright{}, and the root √diś, \textquotedblleft{}to point\textquotedblright{}; MW gives \emph{nir-√diś} \textquotedblleft{}to point to, to show\textquotedblright{}, and \emph{nir-diṣṭa} \textquotedblleft{}pointed out, declared, enjoined\textquotedblright{} (MW).
\end{cousinweb}

\begin{grammarlens}[title={What you've built}]
One more word for giving reasons and opinions.
\end{grammarlens}

\subsection*{Your turn}
\begin{itemize}
\raggedright
  \item \emph{Say it:} \emph{nirdiśati}
  \item \emph{Say it:} \emph{nirdiśati}, once more
  \item \emph{Say it:} \emph{adhyakṣaḥ sadasyān nirdiśati.} — The president directs the members.
  \item \emph{Recall:} say the Sanskrit for foremost, then the Sanskrit for a secretary, then say \emph{nirdiśati} again
\end{itemize}

\begin{tcolorbox}[breakable,colback=teal!4,colframe=teal!35!black,title={Before you move on}]
What does \sk{निर्दिशति} mean? (\textquotedblleft{}Directs; instructs\textquotedblright{}.) Say it once more.
\end{tcolorbox}
